% $MirOS: contrib/code/Flyers/fr.tex,v 1.2 2012/01/27 16:14:40 bsiegert Exp $
%

\documentclass[a4paper,landscape,11pt,notumble]{leaflet}
\usepackage{mirosflyer}

\usepackage[frenchb]{babel}

\hyphenation{Mir-OS}

\begin{document}

%\maketitle
%\thispagestyle{empty}

\graybox
\begin{center}
{\noindent\raisebox{-1ex}{\includegraphics[height=3ex]{img/Flag_of_France}}%
\quad F~~R~~A~~N~~\c C~~A~~I~~S\par}
\vfill

\fontsize{48}{56}\selectfont Le projet\par
\fontsize{64}{64}\selectfont MirOS,\par
\fontsize{48}{56}\selectfont c'est quoi déjà?\par

\vfill

\includegraphics[width=0.75\columnwidth]{img/m-color}

\vfill%
\end{center}%
\graybox

\newpage

\section{Version actuelle}

\subsection{MirOS \#10-current}

\noindent
La version actuellement disponible est MirOS \#10-current, une version bêta de MirOS \#11. Ce système est basé sur OpenBSD-current ainsi que sur des versions antérieures. MirOS BSD inclut également un puissant Shell (mksh), un \emph{ports tree} (MirPorts), une base de donnés d’acronymes (wtf) et GNU cvs 1.12 avec des extensions. 

Le bootloader a été largement réécrit, de façon à être compatible avec le standard multiboot. Il peut donc être appelé par GRUB. 
Le générateur de nombres aléatoires (SRNG) a été amélioré; de nombreux composants du système ont été mis à jour. 
En outre, il est possible d’installer pkgsrc® à la place de MirPorts ou en parallèle.

Le compilateur fourni est gcc 3.4, il intègre le protecteur du stack Propolice et permet la programmation en C, C++, Pascal, Objective-C et Ada.

\subsection{Le DVD}

\noindent
Le DVD distribué lors de cette conférence permet d'installer MirOS~BSD pour les architectures \textcolor{darkred}{i386 et sparc}.
Il n'est pas de Live CD, celui-ci peut être téléchargé sur notre site web.
Pour démarrer l'installation, il suffit de booter à partir du DVD et de sélectionner \flqq~Install\frqq.

Le DVD contient en plus une collection de \textcolor{darkred}{plus de 5000 paquets binaires précompilés} à base d'une version stable de pkgsrc®.
Pour utiliser ces paquets, il faut cocher \url{pkgsrc10.ngz} pendant l'installation du système.
Ensuite, après avoir ajouté \url{/usr/pkg/bin} et \url{/usr/pkg/sbin} au PATH, exécutez la commande:\\
%
\verb,export PKG_PATH=/cdrom/packages/All,\\
%
Maintenant, des paquets binaires peuvent être installés par pkg\_add, par exemple avec ``sudo pkg\_add screen''.

\newpage

\section{MirOS BSD}

\noindent
\largered{MirOS BSD est un système d’exploitation libre et sécurisé de la famille des BSD.}

\noindent
Il est dérivé d’OpenBSD, dont le code source est régulièrement importé et intégré. En outre, il contient du code de NetBSD® et d’autres sources.

MirOS était né suite à un désaccord entre Theo de Raadt, leader d’OpenBSD, et Thorsten Glaser, qui est maintenant le développeur en chef de MirOS. Le responsable pour MirPorts est Benny Siegert. Adam Hoka aide au développement de mksh. Il y a encore d’autres personnes qui contribuent au projet.

MirOS BSD anticipe souvent les grands changements d’OpenBSD. Par exemple, il a été le premier à proposer ELF sur i386 ainsi qu’à intégrer gcc3. Des décisions inverses sont souvent prises par rapport à OpenBSD – par exemple, il n’y aura pas de SMP dans MirOS BSD.

La plupart des différences entre MirOS et OpenBSD se retrouvent dans une infinité de détails – en espérant que le résultat soit un BSD plus convivial.

Le système de base est moins lourd. Des modules peu utilisés (NIS, Kerberos, Bind et BSD Games) ont été supprimés, les deux derniers étant disponibles en tant que port. Le système de base supporte l’internationalisation, UTF-8 et Unicode. Citrus libiconv est également intégré.

Pour les versions stables, des mises à jour binaires sont disponibles lors de l’installation et après.

Le protocole IPv6 fonctionne parfaitement, même dans le serveur web dérivé d’Apache 1.3. Des pilotes pour ISDN sont également fournis.

Pour des raisons legales, nous devons publier les \emph{advertising clauses}. Malheureusement, il y en a trop pour les mentionner ici, mais ils sont tous sur notre site Web
\url{http://www.mirbsd.org/about.htm}.

\newpage

%%% Back Cover Page

\graybox\par
\begin{center}
\vspace{3ex}
{\Large\itshape
Amusez-vous bien avec MirOS!\par\medskip
L'équipe des développeurs\par\medskip
\color{darkred}%
Thorsten Glaser\\
Benny Siegert\par
}

\vfill

\includegraphics[width=0.75\columnwidth]{img/bsdaemon}

\vfill%
\end{center}%
\subsection{Contact}
{\large
\begin{tabular}{ll}
site Web: & \textcolor{darkred}{\url{http://www.mirbsd.org/}}\\
%& http://www.mirbsd.org/\\
IRC: & \textcolor{darkred}{\url{#mirbsd}} sur irc.freenode.net\\
\end{tabular}
\par}
\vspace{3ex}%
\graybox

\newpage

\section{mksh, le Korn Shell de MirOS}

\noindent
L’interpréteur de commande (Shell) par défaut sur MirOS s’appelle mksh – un dérivé du Korn Shell. Les utilisateurs de bash vont s’y retrouver, car la plupart des commandes sont les mêmes.

Cependant, mksh est moins gourmand en mémoire, et plus performant que bash ou zsh. Il unit les caractéristiques de pdksh d’OpenBSD et de Debian ainsi que de nombreux correctifs. Le code est très propre (avertissements, const cleanup); sa sécurité a été contrôlée par les développeurs et les spécialistes de Coverity, Inc.

mksh se prête aussi bien à l’usage interactif qu’à la programmation de scripts. Malgré sa taille modeste, il inclut toutes les fonctionnalités importantes, y compris bien sûr le support pour UTF-8.

mksh est très portable et fonctionne sous de nombreux systèmes, notamment BSD, Linux, Mac OS, Solaris, HURD, HP-UX, AIX et même Windows. Les systèmes suivants (entre autres) fournissent un paquet pour mksh :

\begin{itemize}
\item Debian, Gentoo, FreeWRT, Fedora, d’autres distributions GNU/Linux et OpenSUSE build service
\item NetBSD® pkgsrc®, FreeBSD® Ports, MidnightBSD
\item fink, MacPorts et d’autres collections Mac
\end{itemize}

\section{Pour aller plus loin}

\noindent
MirPorts est notre solution pour installer des logiciels supplémentaires non livrés avec le système de base, nous offrons deux solutions: \textcolor{darkred}{MirPorts} -- dérivé du ports tree d’OpenBSD -- et \textcolor{darkred}{pkgsrc®}, par le projet NetBSD.
Même si tous les paquets ne fonctionnent pas encore sous MirOS~BSD, pkgsrc® contient un plus grand nombre de paquets (plus de 11000) et ceci dans des versions plus récentes.

\newpage

\section{Introduction à MirPorts}

Après son installation ou une mise à jour, la configuration de MirPorts est lancée avec « make setup » dans le dossier \url{/usr/ports}.
Les ports eux-mêmes sont contenus dans des sous-dossiers, classés par catégories.
Un simple « mmake install » dans un dossier télécharge le code source du programme, effectue la compilation, crée un paquet binaire et l’installe.
MirPorts tient compte automatiquement de toutes les dépendences.
Certains ports existent en plusieurs variantes (avec ou sans support pour X, par exemple).

Sur des systèmes MirOS et MirPorts, les \flqq~dotfiles~\frqq, c’est-à-dire les fichiers de configuration cachés, peuvent être rassemblés sous \url{~/.etc}. Ceci évite le chaos habituel dans le dossier home. On peut aussi déposer ses propres logiciels et scripts sous \url{~/.etc/bin}.

MirPorts fonctionne sous les systèmes suivants :
%
\begin{itemize}
\item MirOS BSD (-stable et -current)
\item OpenBSD (-stable et -current)
\item Mac OS X (à partir de 10.4) / Darwin
\item Interix / SFU 3.5 (support partiel)
\end{itemize}

\medskip
\noindent
Même sous les versions stables de MirOS, il est toujours recommandé d’utiliser la version la plus récente de MirPorts. 

\subsection{MirLibtool}

\noindent
Libtool est très utile pour les bibliothèques partagées mais il a des inconvénients. Par exemple, la compilation s’arrête s’il manque le compilateur C++. C’est pourquoi MirPorts contient sa propre version de libtool nommée MirLibtool.

MirLibtool est basé sur GNU libtool 1.5, et est compatible avec tous les versions d’autotools. L’infrastructure de MirPorts l’installe automatiquement  si un port utilise autoconf.

\end{document}
